\documentclass{article}

\usepackage[dblblindworkshop]{neurips_2026}

\workshoptitle{I Can't Believe It's Not Better (ICBINB): Failure Modes of AI in Biology}

\usepackage[utf8]{inputenc}
\usepackage[T1]{fontenc}
\usepackage[hidelinks]{hyperref}
\usepackage{url}
\usepackage{booktabs}
\usepackage{array}
\usepackage{amsfonts}
\usepackage{amsmath}
\usepackage{nicefrac}
\usepackage{microtype}
\usepackage{xcolor}
\usepackage{graphicx}
\usepackage{placeins}

\title{When Protein Language Model Steering Appears to Work: \\
Three Steering Studies and One Endpoint Analysis}

\author{Anonymous Authors}

\begin{document}

\maketitle

\begin{abstract}
Activation steering can change a protein language model's output without
retraining the model, but a higher sequence score does not by itself establish
control of a biological property. We retrospectively examine three completed
ESM2-650M steering studies and one pre-steering endpoint analysis. In the
endpoint analysis, a peptide-only composition score had held-out correlation
$r=0.427$ with a per-peptide label defined as the maximum normalized IC50
score observed across its assayed HLA-II alleles, but only $r=0.100$ with
measured T-cell response. Its full-length correlation changed from $0.379$
under random folds to $-0.323$ under organism-grouped folds. No generation was
run for this target. In the steering studies, an all-layer thermostability
intervention retained 57 to 58 evaluable pairs in a shared low-strength range,
but only 5 and 0 pairs at two higher strengths. An eSol soluble-fraction score
contrast changed from $0.0125$ to $0.00035$ after dominant residues were
excluded, and the resulting interval included zero. A disorder-score contrast
was positive in three stored configurations, but a residue-exclusion interval
excluded zero in only two. These cases motivate this reporting order, but they
do not show that it generalizes beyond the four analyses or that steering
improved any biological property.
\end{abstract}

\section{Problem}
\label{sec:problem}

\subsection{Why apparent steering success needs staged checks}

Protein language models learn statistical structure from large sequence
collections and can support prediction and generation
\citep{lin2023esm}. Activation steering modifies their internal
representations at inference time. A direction is usually estimated from the
difference between activations for high-label and low-label examples, then
added to the residual stream during generation. Prior work applied this
method to protein generation and optimization \citep{huang2025steering}.

The intervention and its evaluation create distinct sources of error. A
generated sequence may receive a higher score because it repeatedly inserts
residues that occur directly in the scoring rule. A conditional score can also
look favorable after low-complexity outputs have been removed from the
comparison. Even before steering, a sequence score may perform well on one
endpoint but poorly on the endpoint named in the claim. These problems concern
the interpretation of the evidence, not only the strength of the intervention.

This paper asks which checks changed the conclusions of three completed
ESM2-650M steering studies and one pre-steering endpoint analysis. We
synthesize the evidence as four retrospective case studies rather than a
catalog of attempted protein properties. The contribution is a proposed
reporting sequence that separates four questions:

\begin{enumerate}
\item Does the score measure the endpoint named in the claim?
\item Do all arms produce comparable outputs, and are the denominators
reported?
\item Does the result remain after dominant residue contributions are
removed?
\item Does the full decision remain consistent across completed run
configurations?
\end{enumerate}

\subsection{Evidence scope}

The synthesis uses three completed steering studies and one completed
pre-steering endpoint analysis. L52 compares a five-layer thermostability
intervention with an all-layer intervention. L56 evaluates
sequence-composition scores against several immune endpoints and stops before
generation. L55 tests intrinsic-disorder steering across three stored
whole-run configurations. L57 tests an eSol soluble-fraction score before and
after excluding dominant residues. The three steering studies use ESM2-650M,
matched-norm random directions, and single-shot masked filling.

The analyses are retrospective. Some checks were added after the original
result had been inspected. We therefore use the studies to show how a check
changes an interpretation, not to estimate the frequency of each failure
mode in protein language model research.

\section{Proposed Approach}
\label{sec:approach}

\subsection{A staged reporting sequence}

Figure~\ref{fig:audit} shows the proposed order of the checks. Each stage
limits the claim that can be carried into the next stage. A future evaluation
should stop at the first unsupported inference. Later checks cannot repair an
earlier failure.

\FloatBarrier
\begin{figure}[htbp]
\centering
\small
\fbox{\parbox[c][2.0cm][c]{0.20\linewidth}{\centering
Endpoint\\
Does the score match the claimed outcome?}}
\hfill
\fbox{\parbox[c][2.0cm][c]{0.20\linewidth}{\centering
Output validity\\
Are outputs comparable and denominators reported?}}
\hfill
\fbox{\parbox[c][2.0cm][c]{0.20\linewidth}{\centering
Composition\\
Does the result survive residue exclusion?}}
\hfill
\fbox{\parbox[c][2.0cm][c]{0.20\linewidth}{\centering
Whole run\\
Does the full decision repeat?}}
\caption{A staged reporting sequence derived from the completed case studies.
A favorable downstream score does not override a failure at an earlier stage.}
\label{fig:audit}
\end{figure}
\FloatBarrier

The endpoint check compares the sequence score with the endpoint named in
the claim. The output check records low-complexity output counts and the number
of source proteins that remain jointly evaluable in every compared arm. The
composition check removes the dominant substituted residues from the score and
repeats the comparison. The whole-run check compares the final decision,
including the composition check, across stored configurations.

\subsection{Case-specific measurements}

L52 evaluated 60 held-out source proteins at five intervention strengths.
Both the five-layer and all-layer interventions used the same per-layer
directions and the same source proteins. The saved analysis compared each
intervention with a matched random direction using 10,000 paired bootstrap
resamples. We inspect the jointly evaluable denominator before interpreting a
conditional score.

L56 used records from the Immune Epitope Database
\citep{vita2019iedb}. Its source table contained 125,985 peptide--allele
binding measurements, covering 15,362 unique peptides and 72 HLA-II alleles.
The analysis grouped the measurements by peptide and used the maximum score
observed across each peptide's assayed alleles. The normalized score was
$1-\log(\mathrm{IC}_{50})/\log(50{,}000)$, so larger values corresponded to
lower IC50 values. The fitted composition score used peptide sequence but not
allele identity. Separate cohorts measured peptide presentation, T-cell
response, and a full-length antigen label. For the full-length cohort, L56
also compared random-fold cross-validation with organism-grouped folds. No
steering intervention was run because the endpoint check failed.

L55 used DisProt annotations \citep{quaglia2022disprot} and the published
TOP-IDP amino-acid scale \citep{campen2008topidp}. L57 used eSol
measurements \citep{niwa2009esol} and an absolute-charge score motivated by
a published recombinant-solubility model
\citep{wilkinson1991solubility}. Each steering study compared the learned
direction with a matched random direction over 150 held-out proteins. The
saved decision used a 95\% paired bootstrap interval over source proteins and
then repeated the score after excluding the two most frequent substituted
residues.

\subsection{Interpretation rules}

We report a conditional score only beside its evaluable denominator. An
interval that includes zero is described as inconclusive, not as evidence of
no effect. A drop under organism-grouped validation is described as
consistent with confounding because grouping does not identify a unique
cause. The three L55 artifacts are treated as completed whole-run
configurations. Their files do not store seed metadata, and one legacy
integer controlled several random processes. Differences among the runs
therefore cannot be attributed to direction construction alone.

\section{Observed Outcome}
\label{sec:outcome}

Table~\ref{tab:cases} summarizes the three failure classes and the corrected
interpretation of each case.

\FloatBarrier
\begin{table}[htbp]
\centering
\caption{Completed evidence used in the retrospective synthesis. The metrics
differ because each check addresses a different inference.}
\label{tab:cases}
\small
\begin{tabular}{@{}>{\raggedright\arraybackslash}p{2.0cm}
                    >{\raggedright\arraybackslash}p{3.2cm}
                    >{\raggedright\arraybackslash}p{4.0cm}
                    >{\raggedright\arraybackslash}p{3.5cm}@{}}
\toprule
Case & Apparent result & Failed or unstable check & Corrected interpretation \\
\midrule
L56 & A peptide-only score predicted an allele-aggregated assay label &
Held-out $r$ was 0.427 for each peptide's maximum normalized IC50 score and
0.100 for T-cell response; organism-grouped $r=-0.323$ &
Validate the endpoint named in the claim and group by source organism \\
\addlinespace
L52 & Five-layer steering appeared favorable at high strength &
All-layer evaluable pairs fell from 57 to 58 in the shared range to 5 and 0 &
Compare methods only where both arms remain evaluable \\
\addlinespace
L55 and L57 & Learned directions increased their sequence scores &
L55 residue exclusion passed in two of three stored runs; the L57 excluded
interval crossed zero &
Report composition and whole-run sensitivity with the main score \\
\bottomrule
\end{tabular}
\end{table}
\FloatBarrier

\subsection{Endpoint choice changed the L56 conclusion}

The fitted peptide-composition score had held-out correlation $r=0.427$ with
a per-peptide label formed by taking the maximum normalized IC50 score across
the alleles assayed for that peptide. The source table spanned 72 HLA-II
alleles, but allele identity was not an input. The result therefore says that
peptide composition was associated with this allele-aggregated dataset label.
It does not estimate affinity for a named HLA-II allele or for a specified
peptide--allele pair. The score's held-out correlation was $r=0.100$ for
measured T-cell response. Fixed motif scores showed the same broad loss of
performance between the binding-assay and T-cell-response cohorts. Performance
on the allele-aggregated assay label therefore did not justify a claim about
immunogenicity.

The full-length cohort exposed a second problem. The composition model had
$r=0.379$ under random five-fold validation. When proteins from the same
organism were kept in the same fold, the correlation was $-0.323$. The mean
within-organism correlation was $0.056$. This pattern is consistent with the
model learning organism-associated composition rather than the intended
endpoint. The evidence does not show that organism is the only cause, and it
does not support a universal claim about sequence-based immune prediction.
It was sufficient to stop this steering target before generation.

\subsection{Low-complexity outputs made the L52 comparison unavailable}

The original L52 decision selected an intervention strength from the full
tested range. All 60 generated outputs retained numerical scores, but the
historical low-complexity filter removed most all-layer outputs at high
strength. At $\alpha=1.0$, the all-layer arm retained only 5 jointly evaluable
pairs. At $\alpha=2.0$, it retained none. The five-layer arm retained 58 and 57
pairs at those strengths. A favorable conditional score for the five-layer arm
therefore did not establish that it was as effective as the all-layer
intervention. The comparison was unavailable under the historical filter.

Restricting the interpretation to the shared low-strength range changed the
result. At $\alpha=0.5$, the five-layer contrast against its random control was
$0.0236$ with a 95\% bootstrap interval of $[0.0202, 0.0270]$. The all-layer
contrast was $0.0498$ with interval $[0.0435, 0.0563]$. Across
$\alpha=0.1$ to $0.5$, the five-layer intervention retained 43\% to 59\% of
the all-layer effect. It produced a score change in this experiment, but the
saved evidence does not support equivalence with all-layer steering.

\subsection{Composition checks changed the L55 and L57 decisions}

At $\alpha=0.5$, the L55 learned-versus-random contrasts were $0.0497$,
$0.0373$, and $0.0435$ in the three stored whole-run configurations. The
historical low-complexity threshold flagged 34, 47, and 33 of 150 learned-arm
outputs, compared with 7, 8, and 13 of 150 random-control outputs. The reported
contrasts used 116, 102, and 113 surviving pairs. The main TOP-IDP contrast was
positive in all three configurations. After E and S were excluded, the
corresponding contrasts were $0.0178$, $0.0039$, and $0.0181$. The bootstrap
interval included zero only in the middle configuration. The positive
conditional contrast repeated, but the residue-exclusion decision did not
repeat in every configuration.

L57 showed a stronger form of composition sensitivity. Its eSol
soluble-fraction score contrast at $\alpha=0.5$ was $0.0125$ with interval
$[0.0086, 0.0166]$. After excluding E and L, the two most frequent
substitutions, the contrast was $0.00035$ with interval
$[-0.0028, 0.0035]$. This does not prove that the remaining effect is
exactly zero. It shows that the original positive decision did not survive
the stated composition check.

The saved one-run geometry result provides limited supporting context. The
L55 and L57 directions had cosine similarity $0.376$ overall and
$0.556$ to $0.666$ in layers 30 to 32. This descriptive overlap has no
control-vector distribution and does not establish a causal relationship
between the two results.

\section{Reason for Failure}
\label{sec:failure}

\subsection{Conditional scores hid denominator loss}

A score computed only on surviving outputs answers a narrow question: how do
the arms differ among source proteins for which both outputs remain
evaluable? It does not describe the intervention over all attempted
generations. In L52, the outputs remained scoreable, but a historical
low-complexity rule left too few all-layer outputs for the high-strength
comparison. Selecting strength by the conditional score allowed the five-layer
method to look favorable where its comparison was unavailable. A steering
report should place output counts, the evaluable denominator, and the
conditional score beside each other, then limit method comparisons to a shared
range.

\subsection{The assay label did not match the claimed endpoint}

The L56 target claim concerned measured T-cell response. The peptide-only
composition score performed better against a per-peptide label defined as the
maximum normalized IC50 score across the alleles assayed for that peptide.
Because the score did not model allele identity, its performance describes an
association with that allele-aggregated label. It is not generic binding
affinity and does not estimate affinity for a specified peptide--allele pair.
It also did not validate T-cell response. Random folds allowed
organism-associated sequence composition to appear in both training and test
sets. The resulting correlation may partly reflect source identity rather than
the target endpoint. Endpoint definitions and grouping rules must therefore be
fixed before a score is used to justify steering.

\subsection{The scoring rule rewarded the dominant substitutions}

Both TOP-IDP and the L57 eSol score are direct functions of amino-acid
composition. Steering can raise such a score by repeatedly inserting residues
that receive favorable weights. Residue exclusion is intentionally strict: it
asks whether the decision survives after the largest direct contribution is
removed. L55 showed that this decision can vary across whole runs even when
the main score change remains positive. L57 showed that a favorable main
interval can change to an interval that includes zero under the same check.

\subsection{Practical reporting procedure}

The cases suggest a concrete order for future evaluations. First, validate
the scoring rule against the endpoint named in the claim and use grouped
splits when source identity may leak. Second, report low-complexity output
counts, evaluable denominators, and conditional scores as separate quantities.
Third, test whether dominant residue substitutions account for the decision.
Finally, repeat the complete decision process across run configurations.
Repeating only the favorable main score is insufficient when an artifact check
carries the conclusion.

\section{Discussion and Limitations}
\label{sec:limitations}

The retrospective synthesis identified different limits because the checks
protect different inferences. L52 concerns denominator loss under a
low-complexity rule. L56 concerns whether a score represents the claimed
endpoint. L55 and L57 concern direct dependence on amino-acid composition and
sensitivity of the final decision. Together, they show why one
successful-looking scalar is not a complete steering evaluation.

The evidence has important limits. Every steering result uses ESM2-650M, one
masked-fill decoder, and one intervention family. The generation studies use
a historical low-complexity threshold and conditional complete-case
comparisons. Their saved files omit exact model revisions, source revisions,
and some run metadata. The three L55 files do not record seed identity, so we
treat them as legacy whole-run configurations. The L56 endpoint analyses use
observational cohorts with different biological contexts and class
structures. The grouped-validation change is consistent with confounding but
does not identify a unique mechanism.

Several checks were selected or clarified after the corresponding outcomes
were known. The bootstrap intervals describe stability over the fixed source
proteins and should not be read as population-wide guarantees. The studies
also contain multiple comparisons without a common familywise correction and
use one matched random direction per saved run. No common audit was applied
prospectively to every case, and no positive controls establish the specificity
of the proposed reporting sequence. Most importantly, every reported endpoint
is computational. No wet-lab assay establishes that a generated sequence has
improved thermostability, intrinsic disorder, soluble fraction,
immunogenicity, or another biological property.

\section{Conclusion}

The completed case studies motivate evaluating protein language model
steering as a sequence of validity checks rather than by one favorable score.
A method comparison can become unavailable when one arm loses most evaluable
outputs. A sequence score may perform differently on the claimed endpoint or
reflect source identity. A positive score can also depend on the residues that
define the score and on the full run configuration. Across these four cases,
reporting the checks together produced a narrower claim that matched the saved
evidence. Whether the same order improves evaluations of other models or
properties remains untested.

\section*{LLM usage disclosure}
OpenAI Codex assisted with manuscript restructuring, prose editing, and source
checks. Earlier project work used Claude (Anthropic) for code and manuscript
drafting. The authors made the scientific decisions and checked the numerical
claims against the saved artifacts.

\section*{Code and data availability}
The study uses public protein and immunology resources together with saved
model outputs. An anonymized artifact package will be provided if the
submission system supports it. The review manuscript does not link to an
identifying repository.

\section*{Ethics declaration}
This computational study used public protein sequences and aggregated public
annotations. It involved no new human participants, clinical intervention, or
collection of personal data.

\section*{Author contributions}
The authors designed the study, conducted the analyses, interpreted the
results, and wrote the manuscript.

\section*{Funding and competing interests}
Funding and competing-interest information is withheld for double-blind
review and will be completed in the final version.

\section*{Reproducibility statement}
The manuscript reports completed artifacts already stored with the project.
The evidence map identifies the source file for every numerical claim and
records known provenance limits. No experiment rerun is part of the remaining
manuscript work.

\bibliographystyle{plainnat}
\bibliography{reference}

\end{document}
